\documentclass[10pt]{article}
\usepackage[paperwidth=612pt,paperheight=792pt,top=34pt,bottom=30pt,left=34pt,right=34pt]{geometry}
\usepackage{amsmath,amssymb}
\usepackage{graphicx}
\usepackage{xcolor}
\usepackage{calc}
\usepackage{fontspec}

% Màu sắc chuẩn theo phong cách sách Stewart Calculus
\definecolor{stewartblue}{RGB}{0, 110, 175}
\definecolor{stewartred}{RGB}{205, 35, 35}
\definecolor{headergray}{RGB}{90, 90, 90}
\definecolor{lightgraytext}{RGB}{120, 120, 120}

\setlength{\parindent}{0pt}
\setlength{\parskip}{0pt}

% Nhãn đánh số phương trình Stewart (khung vuông viền đỏ)
\newcommand{\stewarteqnum}[1]{%
  \begingroup\setlength{\fboxrule}{0.8pt}\setlength{\fboxsep}{2pt}%
  \textcolor{stewartred}{\fbox{\footnotesize\bfseries\textcolor{stewartred}{#1}}}%
  \endgroup%
}

\begin{document}
\pagestyle{empty}

% Tiêu đề đầu trang (Header trang chẵn)
\noindent
\begin{minipage}[b]{136pt}
  {\large\bfseries\textsf{210}}
\end{minipage}%
\hspace{24pt}%
\begin{minipage}[b]{384pt}
  {\bfseries\textsf{\color{headergray}CHƯƠNG 3}}\hspace{1.5em}{\textsf{\color{headergray}Các quy tắc đạo hàm}}
\end{minipage}

\vspace{14pt}

% Khối 1: Phần mở đầu và hai phương trình ban đầu
\noindent
\begin{minipage}[t]{136pt}
  \mbox{}
\end{minipage}%
\hspace{24pt}%
\begin{minipage}[t]{384pt}
  giữa $x$ và $y$, chẳng hạn như
  \vspace{5pt}

  \noindent\stewarteqnum{1}\hfill $x^2 + y^2 = 25$\hfill\mbox{}
  \vspace{4pt}

  \noindent hoặc
  \vspace{4pt}

  \noindent\stewarteqnum{2}\hfill $x^3 + y^3 = 6xy$\hfill\mbox{}
  \vspace{8pt}

  Trong một số trường hợp, ta có thể giải một phương trình như vậy để tìm $y$ dưới dạng một hàm tường minh (hoặc một vài hàm) theo $x$. Chẳng hạn, nếu giải Phương trình 1 theo $y$, ta được $y = \pm\sqrt{25 - x^2}$, vì vậy hai trong số các hàm được xác định bởi Phương trình ẩn 1 là $f(x) = \sqrt{25 - x^2}$ và $g(x) = -\sqrt{25 - x^2}$. Đồ thị của $f$ và $g$ lần lượt là nửa trên và nửa dưới của đường tròn $x^2 + y^2 = 25$. (Xem Hình 1.)
\end{minipage}

\vspace{10pt}

% Khối 2: Hình 1 (Nhãn HÌNH 1 dóng hàng với chân hình ảnh bên phải)
\noindent
\begin{minipage}[b]{136pt}
  {\small\bfseries\textsf{HÌNH 1}}
  \vspace{2pt}
\end{minipage}%
\hspace{24pt}%
\begin{minipage}[b]{384pt}
  \centering
  \includegraphics[width=\linewidth]{images/p0245_fig1.png}
\end{minipage}

\vspace{12pt}

% Khối 3: Đoạn văn về lá Descartes
\noindent
\begin{minipage}[t]{136pt}
  \mbox{}
\end{minipage}%
\hspace{24pt}%
\begin{minipage}[t]{384pt}
  Không dễ để giải Phương trình 2 tìm $y$ một cách tường minh theo $x$ bằng tay. (Ngay cả khi sử dụng công nghệ hỗ trợ, các biểu thức thu được cũng rất phức tạp.) Tuy nhiên, (2) là phương trình của một đường cong được gọi là \textbf{lá Descartes} (folium of Descartes) thể hiện trong Hình 2 và nó xác định $y$ một cách ẩn như một vài hàm số của $x$. Đồ thị của ba hàm số như vậy được thể hiện trong Hình 3. Khi ta nói rằng $f$ là một hàm số được xác định ẩn bởi Phương trình 2, điều đó có nghĩa là phương trình
  \[
    x^3 + [f(x)]^3 = 6x f(x)
  \]
  đúng với mọi giá trị của $x$ thuộc tập xác định của $f$.
\end{minipage}

\vspace{10pt}

% Khối 4: Hình 2 (cột lề trái) và Hình 3 (cột chính)
\noindent
\begin{minipage}[t]{136pt}
  \centering
  \includegraphics[width=\linewidth]{images/p0245_fig2.png}
  \vspace{4pt}
  \raggedright
  {\small\bfseries\textsf{HÌNH 2}}\enspace{\small\textsf{Lá Descartes}}
\end{minipage}%
\hspace{24pt}%
\begin{minipage}[t]{384pt}
  \centering
  \includegraphics[width=\linewidth]{images/p0245_fig3.png}
  \vspace{4pt}
  \raggedright
  {\small\bfseries\textsf{HÌNH 3}}\enspace{\small\textsf{Đồ thị của ba hàm số được xác định ẩn bởi lá Descartes}}
\end{minipage}

\vspace{16pt}

% Khối 5: Mục Đạo hàm ẩn
\noindent
\begin{minipage}[t]{136pt}
  \mbox{}
\end{minipage}%
\hspace{24pt}%
\begin{minipage}[t]{384pt}
  \noindent{\color{stewartblue}\rule{7.5pt}{7.5pt}}\hspace{6pt}{\large\bfseries\textsf{Đạo hàm ẩn}}
  \vspace{6pt}

  May mắn thay, chúng ta không cần phải giải một phương trình để tìm $y$ theo $x$ mới có thể tính được đạo hàm của $y$. Thay vào đó, chúng ta có thể sử dụng phương pháp \textbf{đạo hàm ẩn}. Phương pháp này bao gồm việc lấy đạo hàm cả hai vế của phương trình theo $x$, sau đó giải phương trình thu được để tìm $dy/dx$. Trong các ví dụ và bài tập của mục này, ta luôn giả định rằng phương trình đã cho xác định $y$ một cách ẩn như một hàm khả vi theo $x$ để có thể áp dụng phương pháp đạo hàm ẩn.
\end{minipage}

\vfill

% Dòng bản quyền dưới chân trang
\begin{center}
  \fontsize{6.5pt}{8pt}\selectfont\color{lightgraytext}
  Copyright 2021 Cengage Learning. All Rights Reserved. May not be copied, scanned, or duplicated, in whole or in part. WCN 02-200-203
\end{center}

\end{document}